\section{Problem Setting}

Let $s_t \in \Sstate$ denote a bus-system state at a decision event and $a_t \in \Aact$ denote a holding action in seconds.
The state contains line and station identifiers, time of day, direction, forward and backward headways, waiting passengers, target headway, base dwell, co-line headways, gap, and segment speed.
The transition target is
\[
    \y_t =
    (h^+_{t+1}, h^-_{t+1}, q_{t+1}, d_{t+1}, g_{t+1},
     \bar h^+_{t+1}, \bar h^-_{t+1}, v_{t+1}, r_{t+1}),
\]
covering next headways, waiting passengers, dwell, gap, co-line headways, speed, and reward.

Each city $c$ provides a bundle $\D_c$ of route, timetable, speed, and demand tables.
From these tables we construct paired simulator and static-derived target transitions:
\[
    \y^{\mathrm{sim}} = f_0(s,a), \qquad
    \y^{\mathrm{target}} = f_c(s,a),
\]
and train residual models on
\[
    \Delta_c(s,a) = \y^{\mathrm{target}} - \y^{\mathrm{sim}}.
\]
For cross-city validation, source cities $\mathcal{C}_{\mathrm{src}}$ are used for fitting and a held-out target city $c^\star$ is used for evaluation.
The target transition labels $\y^{\mathrm{target}}$ are never used for target-label-free model fitting or source selection.
However, the target city's unlabeled covariates, schedules, simulator outputs, speed summaries, and static local descriptors may be used to choose or weight source mechanisms.

\paragraph{Target-information budgets.}
We distinguish four information budgets.
In \emph{single-source zero-shot} transfer, one source city provides residual labels and the target city contributes only route/timetable/static-summary inputs.
In \emph{multi-source target-label-free} transfer, several source cities provide residual labels, while the target still provides no transition labels before reporting; unlabeled target summaries may be used for source weighting.
In \emph{few-shot target offline adaptation}, a small passive target history is available before deployment and may update mechanism trust, context bias, or residual sufficient statistics, with held-out target transitions reserved for evaluation.
In \emph{target-calibrated oracle} rows, the target budget is large or in-sample; these rows bound performance but are not no-calibration deployment evidence.
The paper's main claim concerns the second and third budgets, not pure one-city zero-shot deployment.

\paragraph{Traffic-signal stress-test task.}
The RESCO traffic-signal experiment is used as a counterfactual-control stress test rather than as the primary bus-holding task.
For each signalized network $n$, each controlled traffic light $i$, and each decision time $t$, the state $x_{n,i,t}$ contains phase-local queue, occupancy, waiting-time, and downstream-spillback summaries.
The action $a_{i,t}$ is one of the green phases already defined by the SUMO \texttt{tlLogic}; no synthetic phase set is introduced.
The prediction target is the next local queue/cost summary after applying that phase for one control interval.
Cross-network validation holds out one RESCO scenario as the target and fits phase-transfer models from the remaining scenarios.
Zero-target-label rows use only source transitions plus target network descriptors and target-static simulator summaries; few-shot rows may use a short passive target slice to update a low-dimensional bias or trust parameter.
Closed-loop target rollout metrics are used only for final evaluation.

The main evaluation unit is a city target, not a route.
For four cities, the strict leave-one-city-out protocol has exactly four target evaluations:
\[
    \{(\mathcal{C}\setminus\{c\}) \rightarrow c : c \in \mathcal{C}\}.
\]
This avoids counting repeated target cities as independent evidence.
Additional configured split diagnostics are reported separately when they test source-selection behavior rather than the four-city claim.

\paragraph{Claim scope.}
The target transitions are static-data-derived pseudo-observations.
They combine agency data, schedule-derived speeds or headways, passenger-demand summaries, and deterministic transition rules.
They are useful for reproducible cross-city validation, but they are not vehicle-level AVL/APC ground truth.
Accordingly, the paper evaluates transfer robustness under open static-data-derived dynamics, not operational field performance.
Real passive APC/AVL slices validate observed action-$0$ dynamics and target offline adaptation, but they still do not provide labels for unobserved holding actions.
Closed-loop counterfactual policy claims therefore require a validated simulator, benchmark simulator, intervention logs, natural experiments, or field trials; this is why the RESCO traffic-signal benchmark is reported as a counterfactual stress test rather than as bus-field evidence.
